\documentclass{article}
\usepackage{graphicx} % Required for inserting images
\usepackage{hyperref}

\newcommand{\rrs}{\ensuremath{r_{\rm rs}}}
\newcommand{\Rrs}{\ensuremath{R_{\rm rs}}}
\newcommand{\Kd}{\ensuremath{K_{\rm d}}}


% IOPs
\newcommand{\nabs}{\ensuremath{a}} % absorption
\newcommand{\abs}{\ensuremath{a(\lambda)}} % absorption
\newcommand{\scat}{\ensuremath{b(\lambda)}} % back-scattering
\newcommand{\nbb}{\ensuremath{b_b}} % back-scattering
\newcommand{\bb}{\ensuremath{b_b(\lambda)}} % back-scattering
\newcommand{\bbw}{\ensuremath{b_{b,w}(\lambda)}} % back-scattering
\newcommand{\nbbw}{\ensuremath{b_{b,w}}} % back-scattering
\newcommand{\bbnw}{\ensuremath{b_{b,nw}(\lambda)}} % back-scattering
\newcommand{\nbbnw}{\ensuremath{b_{b,nw}}} % back-scattering
\newcommand{\aw}{\ensuremath{a_{\rm w} (\lambda)}}
\newcommand{\naw}{\ensuremath{a_{\rm w}}}
\newcommand{\anw}{\ensuremath{a_{\rm nw} (\lambda)}}
\newcommand{\nanw}{\ensuremath{a_{\rm nw}}}
\newcommand{\acdom}{\ensuremath{a_{\rm CDOM}}}
\newcommand{\aph}{\ensuremath{a_{\rm ph}(\lambda)}}
\newcommand{\naph}{\ensuremath{a_{\rm ph}}}
\newcommand{\adg}{\ensuremath{a_{\rm dg} (\lambda)}}
\newcommand{\nadg}{\ensuremath{a_{\rm dg}}}
\newcommand{\ag}{\ensuremath{a_{\rm g} (\lambda)}}
\newcommand{\ad}{\ensuremath{a_{\rm d} (\lambda)}}
\newcommand{\nad}{\ensuremath{a_{\rm d}}}
\newcommand{\nag}{\ensuremath{a_{\rm g}}}
\newcommand{\ax}{\ensuremath{a_{\rm x}}}
\newcommand{\chl}{\ensuremath{Chl}}
\newcommand{\chla}{\ensuremath{Chla}}

\newcommand{\ffaph}{\ensuremath{\naph(440)}} % aph_L23 at 440
\newcommand{\knwnaph}{\ensuremath{\naph^{\rm L23}(440)}} % aph_L23 at 440

% XIOP
\newcommand{\dl}{\ensuremath{D(\lambda)}}
\newcommand{\avgbbnw}{\ensuremath{<\bbnw>}}
\newcommand{\name}{XQAA}
\newcommand{\claude}{Claude}

% Bricaud

\title{\name\ -- To be renamed}
\author{J. Xavier Prochaska}
\date{July 2024}

\begin{document}

\maketitle

\section{A Simple Derivation}

\subsection{QAA}

As James Hansen first derived decades ago
(1971) and then adopted by Gordon (1973), 
the subsurface remote
sensing reflectance can be approximated
by the quasi single scattering approximation (QSSA):

\begin{equation}
\rrs(\lambda) = G_1 u(\lambda) + G_2 u(\lambda)^2
\label{eqn:hansen}
\end{equation}
with 

\begin{equation}
u(\lambda) \equiv \frac{b_b(\lambda)}{a(\lambda) + b_b(\lambda)}
\label{eqn:def_u}
\end{equation}

\begin{figure}[ht!]
\centering\includegraphics[width=\textwidth]{fig_rrs_u.png}
\caption{
A demonstration that the QSSA is a good descrption
of the full radiative transfer (here the Hydrolight
calculations of Loisel et al. 2023).
}
\label{fig:rrs_u}
\end{figure}

This quantity is shown in Figure~\ref{fig:rrs_u}
for a few wavelengths using the outputs from
Loisel+2023 (hereafter L23).  The community,
following Gordon, has generally assumed that
the coefficients in 
Equation~\ref{eqn:hansen} holds independent
of wavelength.  Figure~\ref{fig:rrs_u} shows
this is not strictly the case, but we shall
also make this assumption for now.  
The formalism that follows can easily accommodate
$G_1, G_2$ as functions of wavelength.

Now, we also can relate the subsurface remote
sensing reflectance to the standard remote
surface reflectance as:

\begin{equation}
    \rrs(\lambda) = \frac{\Rrs(\lambda)}{0.52 + 1.17\Rrs(\lambda)}
\label{eqn:Rrs}
\end{equation}
Clearly, if \rrs\ is held constant, then so too is \Rrs.

\subsection{The Quadratic Equation}

We recognize equation~\ref{eqn:hansen} as the standard,
quadratic equation 

\begin{equation}
    a x^2 + b x + c = 0
\end{equation}
which has solutions:

\begin{equation}
    x = \frac{-b \pm \sqrt{b^2 - 4ac}}{2a}
\end{equation}
Applying this to our equation and taking the positive
root we have:


\begin{equation}
    u = \frac{-G_1 + \sqrt{G_1^2 + 4G_2 \rrs}}{2 G_2}
    \label{eqn:u}
\end{equation}
where $u, \rrs, G_1, G_2$ can all be functions of 
$\lambda$.

Let us define the term on the right-hand side of
Equation~\ref{eqn:u} as a number $C$.
Then, with our definition of $u$ (Equation~\ref{eqn:def_u})
we can re-arrange with a bit of algebra to recover:

\begin{equation}
    \frac{\abs}{\bb} = \frac{1}{C(\lambda)} - 1
    \label{eqn:a_bb}
\end{equation}
Now, let the right-hand side of Equation~\ref{eqn:a_bb}
be another number $D(\lambda) = \equiv 1/C(\lambda) - 1$ and we recover
our \name\ approximation:

\begin{equation}
    \frac{\aw + \anw}{\bbw  + \bbnw} = \dl
    \label{eqn:xiop}
\end{equation}


% %%%%%%%%%%%%%%%%%%%%%%%%%%%%%%%%%%%%%%%%%%%%
\subsection{Red Wavelengths}

At long wavelengths, $\lambda > 600$\,nm, water 
absorption dominates, i.e.\ $\aw \gg \anw$.  
% chl=10; water = 4x phytoplankton
In this
limit, we can solve for the non-water backscattering
coefficients by re-arranging Equation~\ref{eqn:xiop}:

\begin{equation}
    \bbnw \approx \frac{\aw}{\dl} - \bbw
    \label{eqn:bbnw}
\end{equation}
To the extent that we know \bbw\ accurately at these
wavelengths (we do), we have a direct estimate of \bbnw.

% %%%%%%%%%%%%%%%%%%%%%%%%%%%%%%%%%%%%%%%%%%%%
\subsection{Blue Wavelengths}

In most oligotrophic ocean waters ($\chla < 1$), 
  % Equal at \chla=1
  % non-water dominates when \chla=10
backscattering at blue wavelengths
is dominated by pure water: $\bbw \gg \bbnw$.
This allows us to estimate the non-water absorption
coefficients:

\begin{equation}
    \anw \approx \dl \bbw - \aw
    \label{eqn:anw}
\end{equation}

In practice, we find that \bbnw's contribution
is non-negligible.  As such, we may include an estimate of its value from equation~\ref{eqn:bbnw}:

\begin{equation}
    \anw \approx \dl \left [ \bbw + \avgbbnw \right ] - \aw
    \label{eqn:anw_bbnw}
\end{equation}
where \avgbbnw\ is the average 
of \bbnw\ from 600-650\,nm.

% %%%%%%%%%%%%%%%%%%%%%%%%%%%%%%%%%%%%%%%%%%%%%%%%%%%%%%%%%%%%%%%%%%%%%%%%%%%%%%%%
\section{Error Propagation}

Error propagation for \name\ is straightforward and simply follows
standard practice, e.g.

\begin{equation}
    \sigma^2(a) = \left | \frac{\partial a}{\partial \rrs} \right |^2 \sigma^2(\rrs)
\end{equation}

The algebra, however, is a bit hairy..

\subsection{Error in \bbnw}

I have found empirically that \bbnw\ is sensitive to error in $G_1$.
Let's explore why.


\begin{equation}
    \sigma^2(\bbnw) = \left | \frac{\partial \bbnw}{\partial G_1} \right |^2 \sigma^2(G_1)
\end{equation}

Thanks to \claude, we have:

\begin{equation}
    \sigma^2(\bbnw) = \left [ \frac{\aw  + \anw} {(1 - u(\lambda))^2} \right ] 
      * \left [ \frac{-1 + \frac{G_1}{\sqrt{G_1^2 + 4G_2 \rrs}}}{2G_2} \right ]^2 \,  \sigma^2(G_1)
\end{equation}

\end{document}
